\ifdefined\gkver\else\def\gkver{student}\fi
\documentclass[\gkver]{gaokaozhenti}
\begin{document}

\chapter{2022年福建}

\begin{enumerate}
\item
%% number: 1
%% paperName: 2022年福建高考第 1 题 4 分
%% typeId: 1
%% type: 单选题
%% chapter: 功和能
%% point: 功
%% method: 无
%% score: 4
%% degree: 750
%% duplicateId: 0
%% body:
福建土楼兼具居住和防御的功能，承启楼是圆形土楼的典型代表，如图~\ref{2022福建01a}所示。承启楼外楼共四层，各楼层高度如图~\ref{2022福建01b}所示。同一楼层内部通过直径约 $50\Um$ 的圆形廊道连接。若将质量为 $100\Ukg$ 的防御物资先从二楼仓库搬到四楼楼梯口 M 处，再用 $100\Us$ 沿廊道运送到 N 处，如图~\ref{2022福建01c}所示。重力加速度大小取 $10\Umsq$，则\xzanswer{A}

\threepicture[labelA=2022福建01a, labelB=2022福建01b, labelC=2022福建01c, wA=6.0cm, wB=3.0cm, wC=3.0cm]{figs/01a.jpg}{figs/01b.png}{figs/01c.png}

\fourchoices[answer=A]
{该物资从二楼地面被运送到四楼 M 处的过程中，克服重力所做的功为 $5400\UJ$}
{该物资从 M 处被运送到 N 处的过程中，克服重力所做的功为 $78500\UJ$}
{从 M 处沿圆形廊道运动到 N 处，位移大小为 $78.5\Um$}
{从 M 处沿圆形廊道运动到 N 处，平均速率为 $0.5\Ums$}

%% answer: A
%% memo:
\memoanswer{
A．该物资从二楼地面被运送到四楼 M 处的过程中，克服重力所做的功为

$W_{G}=mg\Delta h=100\times10\times(2.7+2.7)\UJ=5400\UJ$

故 A 正确；

B．该物资从 M 处被运送到 N 处的过程中，由于 M、N 高度差为零，所以克服重力做功为零，故 B 错误；

C．从 M 处沿圆形廊道运动到 N 处，位移大小为 $50\Um$，故 C 错误；

D．从 M 处沿圆形廊道运动到 N 处，平均速率为

$v=\frac{s}{t}=\frac{\pi r}{t}=\frac{3.14\times\frac{50}{2}}{100}\Ums=0.785\Ums$

故 D 错误。

故选 A。
}

\item
%% number: 2
%% paperName: 2022年福建高考第 2 题 4 分
%% typeId: 1
%% type: 单选题
%% chapter: 原子和原子核
%% point: 半衰期
%% method: 无
%% score: 4
%% degree: 650
%% duplicateId: 0
%% body:
2011 年 3 月，日本发生的大地震造成了福岛核电站核泄漏。在泄露的污染物中含有大量放射性元素 \ce{^{131}_{53}I}，其衰变方程为 \ce{^{131}_{53}I -> ^{131}_{54}Xe + ^{0}_{-1}e}，半衰期为 8 天，已知 $m_{I}=131.037\,21\Uu$，$m_{Xe}=131.031\,86\Uu$，$m_{e}=0.000\,549\Uu$，则下列说法正确的是\xzanswer{D}

\fourchoices[answer=D]
{衰变产生的 $\beta$ 射线来自于 \ce{^{131}_{53}I} 原子的核外电子}
{该反应前后质量亏损 $0.005\,35\Uu$}
{放射性元素 \ce{^{131}_{53}I} 发生的衰变为 $\alpha$ 衰变}
{经过 16 天，$75\%$ 的 \ce{^{131}_{53}I} 原子核发生了衰变}

%% answer: D
%% memo:
\memoanswer{
A．\ce{^{131}_{53}I} 衰变时，原子核内中子转化为质子和电子，大量电子从原子核释放出来形成 $\beta$ 射线，故 A 错误；

B．该反应前后质量亏损为

$\Delta m=m_{I}-m_{Xe}-m_{e}=131.037\,21\Uu-131.031\,86\Uu-0.000\,549\Uu=0.004\,801\Uu$

故 B 错误；

C．放射性元素 \ce{^{131}_{53}I} 发生的衰变为 $\beta$ 衰变，故 C 错误；

D．由于半衰期为 8 天，可知经过 16 天，即经过两个半衰期，$75\%$ 的 \ce{^{131}_{53}I} 原子核发生了衰变，故 D 正确。

故选 D。
}

\item
%% number: 3
%% paperName: 2022年福建高考第 3 题 4 分
%% typeId: 1
%% type: 单选题
%% chapter: 电场
%% point: 电场线
%% method: 无
%% score: 4
%% degree: 650
%% duplicateId: 0
%% body:
平时我们所处的地球表面，实际上存在场强大小为 $100\UVm$ 的电场，可将其视为匀强电场，在地面立一金属杆后空间中的等势面如图所示。空间中存在 a、b、c 三点，其中 a 点位于金属杆正上方，b、c 等高。则下列说法正确的是\xzanswer{B}

\onepicture[w=5.0cm, align=right]{figs/03.png}

\fourchoices[answer=B]
{b、c 两点的电势差 $U_{bc}=0$}
{a 点场强大小大于 $100\UVm$}
{a 点场强方向水平向右}
{a 点的电势低于 c 点}

%% answer: B
%% memo:
\memoanswer{
A．由图可知，b、c 两点的电势差为 $U_{bc}=(200-300)\UV=-100\UV$。故 A 错误；

B．由图可知，a 点与相邻两等势面的距离小于 $1\Um$，电势差等于 $100\UV$，根据 $E=\frac{U}{d}$ 可知 a 点场强大小大于 $100\UVm$，故 B 正确；

C．根据场强方向垂直于等势面，可知 a 点的场强方向沿竖直方向，不是水平方向，故 C 错误；

D．由图可知，a 点与 c 点在同一等势面上，电势均为 $300\UV$，故 D 错误。

故选 B。
}

\item
%% number: 4
%% paperName: 2022年福建高考第 4 题 4 分
%% typeId: 1
%% type: 单选题
%% chapter: 曲线运动
%% point: 人造地球卫星
%% method: 无
%% score: 4
%% degree: 450
%% duplicateId: 0
%% body:
2021 年美国“星链”卫星曾近距离接近我国运行在距地 $390\Ukm$ 近圆轨道上的天宫空间站。为避免发生危险，天宫空间站实施了发动机点火变轨的紧急避碰措施。已知质量为 $m$ 的物体从距地心 $r$ 处运动到无穷远处克服地球引力所做的功为 $G\frac{Mm}{r}$，式中 $M$ 为地球质量，$G$ 为引力常量；现将空间站的质量记为 $m_{0}$，变轨前后稳定运行的轨道半径分别记为 $r_{1}$、$r_{2}$，如图所示。空间站紧急避碰过程发动机做的功至少为\xzanswer{A}

\onepicture[w=5.0cm, align=right]{figs/04.png}

\fourchoices[answer=A]
{$\frac{1}{2}GMm_{0}\left(\frac{1}{r_{1}}-\frac{1}{r_{2}}\right)$}
{$GMm_{0}\left(\frac{1}{r_{1}}-\frac{1}{r_{2}}\right)$}
{$\frac{3}{2}GMm_{0}\left(\frac{1}{r_{1}}-\frac{1}{r_{2}}\right)$}
{$2GMm_{0}\left(\frac{1}{r_{1}}-\frac{1}{r_{2}}\right)$}

%% answer: A
%% memo:
\memoanswer{
空间站紧急避碰的过程可简化为加速、变轨、再加速的三个阶段；空间站从轨道 $r_{1}$ 变轨到 $r_{2}$ 过程，根据动能定理有

$W+W_{\text{引力}}=\Delta E_{k}$

依题意可得引力做功

$W_{\text{引力}}=G\frac{Mm_{0}}{r_{2}}-G\frac{Mm_{0}}{r_{1}}$

万有引力提供在圆形轨道上做匀速圆周运动的向心力，由牛顿第二定律有

$G\frac{Mm_{0}}{r^{2}}=m_{0}\frac{v^{2}}{r}$

求得空间站在轨道上运动的动能为

$E_{k}=G\frac{Mm_{0}}{2r}$

动能的变化

$\Delta E_{k}=G\frac{Mm_{0}}{2r_{2}}-G\frac{Mm_{0}}{2r_{1}}$

解得 $W=\frac{1}{2}GMm_{0}\left(\frac{1}{r_{1}}-\frac{1}{r_{2}}\right)$

故选 A。
}

\item
%% number: 5
%% paperName: 2022年福建高考第 5 题 6 分
%% typeId: 2
%% type: 多选题
%% chapter: 磁场
%% point: 电流的磁场
%% method: 无
%% score: 6
%% degree: 650
%% duplicateId: 0
%% body:
奥斯特利用如图所示实验装置研究电流的磁效应。一个可自由转动的小磁针放在白金丝导线正下方，导线两端与一伏打电池相连。接通电源瞬间，小磁针发生了明显偏转。奥斯特采用控制变量法，继续研究了导线直径、导线材料、电池电动势以及小磁针位置等因素对小磁针偏转情况的影响。他能得到的实验结果有\xzanswer{AB}

\onepicture[w=5.0cm, align=right]{figs/05.png}

\fourchoices[answer=AB]
{减小白金丝直径，小磁针仍能偏转}
{用铜导线替换白金丝，小磁针仍能偏转}
{减小电源电动势，小磁针一定不能偏转}
{小磁针的偏转情况与其放置位置无关}

%% answer: AB
%% memo:
\memoanswer{
A．减小导线直径，仍存在电流，其产生的磁场仍能使小磁针偏转，选项 A 正确；

B．白金导线换成铜导线，仍存在电流，产生的磁场仍能使小磁针偏转，选项 B 正确；

C．减小伏打电池电动势，只要导线中有电流，小磁场还是会发生偏转，选项 C 错误；

D．通电导线产生的磁场与地磁场叠加后，其空间磁场方向与位置有关，当小磁针在不同位置时其偏转情况不同，选项 D 错误。

故选 AB。
}

\item
%% number: 6
%% paperName: 2022年福建高考第 6 题 6 分
%% typeId: 2
%% type: 多选题
%% chapter: 交流电
%% point: 变压器
%% method: 无
%% score: 6
%% degree: 550
%% duplicateId: 0
%% body:
某同学利用如图所示电路模拟远距离输电．图中交流电源电压为 $6\UV$，定值电阻 $R_{1}=R_{2}=20\UO$，小灯泡 L$_{1}$、L$_{2}$ 的规格均为“$6\UV$  $1.8\UW$”，理想变压器 T$_{1}$、T$_{2}$ 原副线圈的匝数比分别为 $1:3$ 和 $3:1$．分别接通电路 Ⅰ 和电路 Ⅱ，两电路都稳定工作时，\xzanswer{BC}

\onepicture[w=6.2cm, align=right]{figs/06.jpg}

\fourchoices[answer=BC]
{L$_{1}$ 与 L$_{2}$ 一样亮}
{L$_{2}$ 比 L$_{1}$ 更亮}
{$R_{1}$ 上消耗的功率比 $R_{2}$ 的大}
{$R_{1}$ 上消耗的功率比 $R_{2}$ 的小}

%% answer: BC
%% memo:
\memoanswer{
若开关接 cd 端，则若电源电压为 $U_{0}$，理想变压器 $T_{1}$、$T_{2}$ 的匝数比为

$\frac{n_{2}}{n_{1}}=\frac{n_{3}}{n_{4}}=k$

用户电阻为 $R_{\text{负载}}$，输电线电阻为 $R_{\text{导线}}$，由变压器工作原理和欧姆定律。升压变压器次级电压

$U_{2}=kU_{0}$

降压变压器初级电压

$U_{3}=U_{2}-I_{2}R_{\text{导线}}$

降压变压器次级电压

$U_{4}=\frac{U_{3}}{k}$

$I_{4}=\frac{U_{4}}{R_{\text{负载}}}$

$\frac{I_{4}}{I_{3}}=k$

$I_{3}=I_{2}$

可得输电功率为

$P_{\text{输}}=U_{2}I_{2}=\frac{k^{2}U_{0}^{2}}{R_{\text{导线}}+k^{2}R_{\text{负载}}}$

输电线上损耗的电功率为

$P_{\text{导线}}=I_{2}^{2}R_{\text{导线}}=\frac{k^{2}U_{0}^{2}}{\left(R_{\text{导线}}+k^{2}R_{\text{负载}}\right)^{2}}R_{\text{导线}}$

用户得到的电功率为

$P_{\text{负载}}=\frac{k^{2}U_{0}^{2}}{\left(R_{\text{导线}}+k^{2}R_{\text{负载}}\right)^{2}}\cdot k^{2}R_{\text{负载}}$

若开关接 ab 端，则负载得到的功率

$P_{\text{负载}}^{\prime}=\frac{U_{0}^{2}}{\left(R_{\text{导线}}+R_{\text{负载}}\right)^{2}}\cdot R_{\text{负载}}$

输电线上损耗的电功率为

$P_{\text{导线}}^{\prime}=\frac{U_{0}^{2}}{\left(R_{\text{导线}}+R_{\text{负载}}\right)^{2}}R_{\text{导线}}$

将 $R_{\text{导线}}=R_{1}=R_{2}=20\UO$，$R_{\text{负载}}=\frac{6^{2}}{1.8}\UO=20\UO$，$k=3$ 带入可知：

可得 $P_{\text{负载}} > P_{\text{负载}}^{\prime}$

即 $L_{2}$ 比 $L_{1}$ 更亮；

$P_{\text{导线}} < P_{\text{导线}}^{\prime}$

$R_{1}$ 上消耗的功率比 $R_{2}$ 的大。

故选 BC。
}

\item
%% number: 7
%% paperName: 2022年福建高考第 7 题 6 分
%% typeId: 2
%% type: 多选题
%% chapter: 功和能
%% point: 功能中的图象
%% method: 无
%% score: 6
%% degree: 450
%% duplicateId: 0
%% body:
一物块以初速度 $v_{0}$ 自固定斜面底端沿斜面向上运动，一段时间后回到斜面底端。该物体的动能 $E_{k}$ 随位移 $x$ 的变化关系如图所示，图中 $x_{0}$、$E_{k1}$、$E_{k2}$ 均己知。根据图中信息可以求出的物理量有\xzanswer{BD}

\onepicture[w=5.0cm, align=right]{figs/07.png}

\fourchoices[answer=BD]
{重力加速度大小}
{物体所受滑动摩擦力的大小}
{斜面的倾角}
{沿斜面上滑的时间}

%% answer: BD
%% memo:
\memoanswer{
ABC．由动能定义式得 $E_{k1}=\frac{1}{2}mv_{0}^{2}$，则可求解质量 $m$；上滑时，由动能定理 $E_{k}-E_{k1}=-(mg\sin\theta+f)x$

下滑时，由动能定理

$E_{k}=(mg\sin\theta-f)(x_{0}-x)$

$x_{0}$ 为上滑的最远距离；由图像的斜率可知

$mg\sin\theta+f=\frac{E_{k1}}{x_{0}}$，$mg\sin\theta-f=\frac{E_{k2}}{x_{0}}$

两式相加可得

$g\sin\theta=\frac{1}{2m}\left(\frac{E_{k1}}{x_{0}}+\frac{E_{k2}}{x_{0}}\right)$

相减可知

$f=\frac{E_{k1}-E_{k2}}{2x_{0}}$

即可求解 $g\sin\theta$ 和所受滑动摩擦力 $f$ 的大小，但重力加速度大小、斜面的倾角不能求出，故 AC 错误，B 正确；

D．根据牛顿第二定律和运动学关系得

$mg\sin\theta+f=ma$，$t=\frac{v_{0}}{a}$

故可求解沿斜面上滑的时间，D 正确。

故选 BD。
}

\item
%% number: 8
%% paperName: 2022年福建高考第 8 题 6 分
%% typeId: 2
%% type: 多选题
%% chapter: 运动和力
%% point: 动量定理
%% method: 无
%% score: 6
%% degree: 450
%% duplicateId: 0
%% body:
我国霍尔推进器技术世界领先，其简化的工作原理如图所示。放电通道两端电极间存在一加速电场，该区域内有一与电场近似垂直的约束磁场（未画出）用于提高工作物质被电离的比例。工作时，工作物质氙气进入放电通道后被电离为氙离子，再经电场加速喷出，形成推力。某次测试中，氙气被电离的比例为 $95\%$，氙离子喷射速度为 $1.6\times10^{4}\Ums$，推进器产生的推力为 $80\UmN$。已知氙离子的比荷为 $7.3\times10^{5}\UCkg$；计算时，取氙离子的初速度为零，忽略磁场对离子的作用力及粒子之间的相互作用，则\xzanswer{AD}

\onepicture[w=6.2cm, align=right]{figs/08.png}

\fourchoices[answer=AD]
{氙离子的加速电压约为 $175\UV$}
{氙离子的加速电压约为 $700\UV$}
{氙离子向外喷射形成的电流约为 $37\UA$}
{每秒进入放电通道的氙气质量约为 $5.3\times10^{-6}\Ukg$}

%% answer: AD
%% memo:
\memoanswer{
AB．氙离子经电场加速，根据动能定理有

$qU=\frac{1}{2}mv^{2}-0$

可得加速电压为

$U=\frac{v^{2}}{2\left(\frac{q}{m}\right)}\approx175\UV$

故 A 正确，B 错误；

D．在 $\Delta t$ 时间内，有质量为 $\Delta m$ 的氙离子以速度 $v$ 喷射而出，形成电流为 $I$，由动量定理可得

$F\Delta t=\Delta mv-0$

进入放电通道的氙气质量为 $\Delta m_{0}$，被电离的比例为 $\eta$，则有

$\frac{\Delta m}{\Delta t}=\eta\left(\frac{\Delta m_{0}}{\Delta t}\right)$

联立解得

$\frac{\Delta m_{0}}{\Delta t}=\frac{F}{\eta v}\approx5.3\times10^{-6}\Ukg$

故 D 正确；

C．在 $\Delta t$ 时间内，有电荷量为 $\Delta Q$ 的氙离子喷射出，则有

$\Delta Q=\left(\frac{\Delta m}{m}\right)q$，$I=\frac{\Delta Q}{\Delta t}$

联立解得

$I=\left(\frac{F}{v}\right)\left(\frac{q}{m}\right)\approx3.7\UA$

故 C 错误。

故选 AD。
}

\item
%% number: 9
%% paperName: 2022年福建高考第 9 题 4 分
%% typeId: 3
%% type: 填空题
%% chapter: 机械波
%% point: 波长频率波速
%% method: 无
%% score: 4
%% degree: 650
%% duplicateId: 0
%% body:
艺术体操运动员站在场地中以一定频率上下抖动 $6\Um$ 长绸带的一端，绸带自左向右呈现波浪状起伏。某时刻绸带形状如图所示（符合正弦函数特征），此时绸带上 P 点运动方向 \tkanswer[1.4cm]{向上}（填“向上”、“向下”、“向左”或“向右”）。保持抖动幅度不变，如果要在该绸带上产生更加密集的波浪状起伏效果，运动员上下抖动的频率应 \tkanswer[1.4cm]{增大}（填“增大”、“减小”或“保持不变”）。

\onepicture[w=6.0cm, align=right]{figs/09.png}

%% answer:
%% memo:
\memoanswer{
（1）从图中可知绸带上形成的波是自左向右传播的，根据波形平移法，可判断绸带上 P 点运动方向向上；

（2）绸带上产生更加密集的波浪状起伏效果，说明波长变小，而同种介质中同类型波的传播波速是不变的，根据 $f=\frac{v}{\lambda}$ 可知运动员上下抖动的频率增大。
}

\item
%% number: 10
%% paperName: 2022年福建高考第 10 题 4 分
%% typeId: 3
%% type: 填空题
%% chapter: 气体性质
%% point: 气体图象
%% method: 无
%% score: 4
%% degree: 650
%% duplicateId: 0
%% body:
带有活塞的汽缸内封闭一定质量的理想气体，气体开始处于 a 状态，然后经过 a$\to$b$\to$c 状态变化过程到达 c 状态。在 $V-T$ 图中变化过程如图所示。

（1）气体从 a 状态经过 a$\to$b 到达 b 状态的过程中压强 \tkanswer[1.6cm]{增大}。（填“增大”、“减小”或“不变”）

（2）气体从 b 状态经过 b$\to$c 到达 c 状态的过程要 \tkanswer[1.4cm]{放出}（填“吸收”或“放出”）热量。

\onepicture[w=5.0cm, align=right]{figs/10.png}

%% answer:
%% memo:
\memoanswer{
（1）由 $V-T$ 图像可知，气体从 a 状态经过 a$\to$b 到达 b 状态的过程中，气体的体积保持不变，温度升高，根据查理定律可知气体的压强增大。

（2）由 $V-T$ 图像可知，气体从 $b$ 状态经过 b$\to$c 到达 c 状态的过程，气体的温度保持不变，则气体的内能保持不变；气体的体积减小，则外界对气体做正功，根据热力学第一定律可知，气体对外放出热量。
}

\item
%% number: 11
%% paperName: 2022年福建高考第 11 题 5 分
%% typeId: 4
%% type: 实验题
%% chapter: 曲线运动
%% point: 研究平抛运动实验
%% method: 无
%% score: 5
%% degree: 550
%% duplicateId: 0
%% body:
某实验小组利用图~\ref{2022福建11a}所示装置验证小球平抛运动的特点。实验时，先将斜槽固定在贴有复写纸和白纸的木板边缘，调节槽口水平并使木板竖直；把小球放在槽口处，用铅笔记下小球在槽口时球心在木板上的水平投影点 $O$，建立 $xOy$ 坐标系．然后从斜槽上固定的位置释放小球，小球落到挡板上并在白纸上留下印迹．上下调节挡板进行多次实验．实验结束后，测量各印迹中心点 $O_{1}$、$O_{2}$、$O_{3} \ldots$ 的坐标，并填入表格中，计算对应的 $x^{2}$ 值。

\begin{center}
\begin{tblr}{|c|c|c|c|c|c|c|}
\hline
 & $O_{1}$ & $O_{2}$ & $O_{3}$ & $O_{4}$ & $O_{5}$ & $O_{6}$ \\
\hline
$y/\Ucm$ & 2.95 & 6.52 & 9.27 & 13.20 & 16.61 & 19.90 \\
\hline
$x/\Ucm$ & 5.95 & 8.81 & 10.74 & 12.49 & 14.05 & 15.28 \\
\hline
$x^{2}/\Ucmq$ & 35.4 & 77.6 & 115.3 & 156.0 & 197.4 & 233.5 \\
\hline
\end{tblr}
\end{center}

\begin{enumerate}
	\item
	根据上表数据，在图~\ref{2022福建11b}给出的坐标纸上补上 $O_{4}$ 数据点，并绘制“$y-x^{2}$”图线。
	\item
	由 $y-x^{2}$ 图线可知，小球下落的高度 $y$，与水平距离的平方 $x^{2}$ 成 \tkanswer[1.4cm]{线性}（填“线性”或“非线性”）关系，由此判断小球下落的轨迹是抛物线。
	\item
	由 $y-x^{2}$ 图线求得斜率 $k$，小球平抛运动的初速度表达式为 $v_{0}=$ \tkanswer[2.2cm]{$\sqrt{\frac{g}{2k}}$}（用斜率 $k$ 和重力加速度 $g$ 表示）。
	\item
	该实验得到的 $y-x^{2}$ 图线常不经过原点，可能的原因是 \tkanswer[5.0cm]{水平射出点未与 O 点重合}。
\end{enumerate}

\twopicture[labelA=2022福建11a, labelB=2022福建11b, wA=3.6cm, wB=6.0cm, align=right]{figs/11a.png}{figs/11b.png}

%% answer:
\jdanswer{
\begin{enumerate}
	\item
	如图所示
	\item
	线性
	\item
	$\sqrt{\frac{g}{2k}}$
	\item
	水平射出点未与 $O$ 点重合
\end{enumerate}
}
%% memo:
\memoanswer{
（1）根据上表数据在坐标纸上描出 $O_{4}$ 数据点，并绘制“$y-x^{2}$”图线如图所示

\onepicture[w=5.5cm]{figs/11_an.jpg}

（2）由 $y-x^{2}$ 图线为一条倾斜的直线可知，小球下落的高度 $y$，与水平距离的平方 $x^{2}$ 成线性关系。

（3）根据平抛运动规律可得

$x=v_{0}t$

$y=\frac{1}{2}gt^{2}$

联立可得

$y=\frac{1}{2}g\left(\frac{x}{v_{0}}\right)^{2}=\frac{g}{2v_{0}^{2}}x^{2}$

可知 $y-x^{2}$ 图像的斜率为

$k=\frac{g}{2v_{0}^{2}}$

解得小球平抛运动的初速度为

$v_{0}=\sqrt{\frac{g}{2k}}$

（4）$y-x^{2}$ 图线是一条直线，但常不经过原点，说明实验中测量的 $y$ 值偏大或偏小一个定值，这是小球的水平射出点未与 $O$ 点重合，位于坐标原点 $O$ 上方或下方所造成的。
}

\item
%% number: 12
%% paperName: 2022年福建高考第 12 题 7 分
%% typeId: 4
%% type: 实验题
%% chapter: 电路
%% point: 测电动势与内阻
%% method: 无
%% score: 7
%% degree: 450
%% duplicateId: 0
%% body:
在测量某电源电动势和内阻时，因为电压表和电流表的影响，不论使用何种接法，都会产生系统误差，为了消除电表内阻造成的系统误差，某实验兴趣小组设计了如图~\ref{2022福建12a}实验电路进行测量。已知 $R_{0}=2\UO$。

\begin{enumerate}
	\item
	按照图~\ref{2022福建12a}所示的电路图，将图~\ref{2022福建12b}中的器材实物连线补充完整。
	\item
	实验操作步骤如下：
	\begin{steps}
		\item
		将滑动变阻器滑到最左端位置
		\item
		接法 Ⅰ：单刀双掷开关 S 与 1 接通，闭合开关 S$_{0}$，调节滑动变阻器 $R$，记录下若干组数据 $U_{1}-I_{1}$ 的值，断开开关 S$_{0}$
		\item
		将滑动变阻器滑到最左端位置
		\item
		接法 Ⅱ：单刀双掷开关 S 与 2 闭合，闭合开关 S$_{0}$，调节滑动变阻器 $R$，记录下若干组数据 $U_{2}-I_{2}$ 的值，断开开关 S$_{0}$
		\item
		分别作出两种情况所对应的 $U_{1}-I_{1}$ 和 $U_{2}-I_{2}$ 图像
	\end{steps}
	\item
	单刀双掷开关接 1 时，某次读取电表数据时，电压表指针如图~\ref{2022福建12c}所示，此时 $U_{1}=$ \tkanswer[1.2cm]{1.30} $\UV$。
	\item
	根据测得数据，作出 $U_{1}-I_{1}$ 和 $U_{2}-I_{2}$ 图像如图~\ref{2022福建12d}所示，根据图线求得电源电动势 $E=$ \tkanswer[1.2cm]{1.80} $\UV$，内阻 $r=$ \tkanswer[1.2cm]{2.50} $\UO$。（结果均保留两位小数）
	\item
	由图~\ref{2022福建12d}可知 \tkanswer[1.6cm]{接法 Ⅱ}（填“接法 Ⅰ”或“接法 Ⅱ”）测得的电源内阻更接近真实值。
	\item
	综合考虑，若只能选择一种接法，应选择 \tkanswer[1.6cm]{接法 Ⅱ}（填“接法 Ⅰ”或“接法 Ⅱ”）测量更合适。
\end{enumerate}

\fourpicture[labelA=2022福建12a, labelB=2022福建12b, labelC=2022福建12c, labelD=2022福建12d, numA=甲, numB=乙, numC=丙, numD=丁, wA=4.6cm, wB=6.4cm, wC=3.4cm, wD=6.4cm, align=right]{figs/12a.png}{figs/12b.png}{figs/12c.jpg}{figs/12d.jpg}

%% answer:
\jdanswer{
\begin{enumerate}
	\item
	如图所示
	\item

	\item
	1.30
	\item
	1.80，2.50
	\item
	接法 Ⅱ
	\item
	接法 Ⅱ
\end{enumerate}
}
%% memo:
\memoanswer{
（1）根据图甲所示的电路图，实物连接如图所示

\onepicture[w=6.0cm]{figs/12_an.jpg}

（3）量程为 $3\UV$ 的电压表分度值为 $0.1\UV$，需要估读到分度值的下一位，由图丙可知电压表读数为 $U_{1}=1.30\UV$

（4）当单刀双掷开关接 1 时，电流表示数为零时，电压表测量准确，故电动势为 $U_{1}-I_{1}$ 的纵轴截距，则有 $E=1.80\UV$

当单刀双掷开关接 2 时，电压表示数为零时，电流表测量准确，由 $U_{2}-I_{2}$ 图像可知此时电路电流为 $0.40\UA$，根据闭合电路欧姆定律可知

$I=\frac{E}{R_{0}+r}$

解得内阻为

$r=\frac{E}{I}-R_{0}=\frac{1.80}{0.40}\UO-2\UO=2.50\UO$

（5）由图丁可知 $U_{1}-I_{1}$ 图像的斜率为

$k_{1}=\frac{1.80-0}{0.36}\UO=R_{0}+r_{1}$

解得 $r_{1}=3.00\UO$

由图丁可知 $U_{2}-I_{2}$ 图像的斜率为

$k_{2}=\frac{1.70-0}{0.40}\UO=R_{0}+r_{2}$

解得 $r_{2}=2.25\UO$

可得 $\frac{r_{1}-r}{r}=\frac{3.00-2.50}{2.50}=0.2>\frac{r-r_{2}}{r}=\frac{2.50-2.25}{2.50}=0.1$

故接法Ⅱ测得的电源内阻更接近真实值。

（6）由电路图可知接法Ⅰ的误差来源是电流表的分压，接法Ⅱ的误差来源是电压表的分流，由于电源内阻较小，远小于电压表内阻，结合（5）问分析可知，若只能选择一种接法，应选择接法Ⅱ更合适。
}

\item
%% number: 13
%% paperName: 2022年福建高考第 13 题 12 分
%% typeId: 6
%% type: 计算题
%% chapter: 曲线运动
%% point: 圆周运动的应用
%% method: 无
%% score: 12
%% degree: 650
%% duplicateId: 0
%% body:
清代乾隆的《冰嬉赋》用“躄躠”（可理解为低身斜体）二字揭示了滑冰的动作要领。$500\Um$ 短道速滑世界纪录由我国运动员武大靖创造并保持。在其创造纪录的比赛中，

\begin{enumerate}
	\item
	武大靖从静止出发，先沿直道加速滑行，前 $8\Um$ 用时 $2\Us$。该过程可视为匀加速直线运动，求此过程加速度大小；
	\item
	武大靖途中某次过弯时的运动可视为半径为 $10\Um$ 的匀速圆周运动，速度大小为 $14\Ums$。已知武大靖的质量为 $73\Ukg$，求此次过弯时所需的向心力大小；
	\item
	武大靖通过侧身来调整身体与水平冰面的夹角，使场地对其作用力指向身体重心而实现平稳过弯，如图所示。求武大靖在（2）问中过弯时身体与水平面的夹角 $\theta$ 的大小。（不计空气阻力，重力加速度大小取 $10\Umsq$，$\tan 22^\circ=0.40$、$\tan 27^\circ=0.51$、$\tan 32^\circ=0.62$、$\tan 37^\circ=0.75$）
\end{enumerate}

\onepicture[w=6.0cm, align=right]{figs/13.png}

%% answer:
\jdanswer{
\begin{enumerate}
	\item
	$4\Umsq$
	\item
	$1430.8\UN$
	\item
	$27^\circ$
\end{enumerate}
}
%% memo:
\memoanswer{
（1）设武大靖运动过程的加速度大小为 $a$，根据 $x=\frac{1}{2}at^{2}$

解得 $a=\frac{2x}{t^{2}}=\frac{2\times8}{2^{2}}\Umsq=4\Umsq$

（2）根据 $F_{\text{向}}=m\frac{v^{2}}{r}$

解得过弯时所需的向心力大小为

$F_{\text{向}}=73\times\frac{14^{2}}{10}\UN=1430.8\UN$

（3）设场地对武大靖的作用力大小为 $F$，受力如图所示

根据牛顿第二定律可得 $F_{\text{向}}=\frac{mg}{\tan\theta}$

\onepicture[w=2.6cm]{figs/13_an.jpg}

$\tan\theta=\frac{mg}{F_{\text{向}}}=\frac{73\times10}{1430.8}\approx0.51$

解得

可得 $\theta=27^\circ$
}

\item
%% number: 14
%% paperName: 2022年福建高考第 14 题 12 分
%% typeId: 6
%% type: 计算题
%% chapter: 运动和力
%% point: 动量守恒定律
%% method: 无
%% score: 12
%% degree: 550
%% duplicateId: 0
%% body:
如图，L 形滑板 A 静置在粗糙水平面上，滑板右端固定一劲度系数为 $k$ 的轻质弹簧，弹簧左端与一小物块 B 相连，弹簧处于原长状态。一小物块 C 以初速度 $v_{0}$ 从滑板最左端滑入，滑行 $s_{0}$ 后与 B 发生完全非弹性碰撞（碰撞时间极短），然后一起向右运动；一段时间后，滑板 A 也开始运动．已知 A、B、C 的质量均为 $m$，滑板与小物块、滑板与地面之间的动摩擦因数均为 $\mu$，重力加速度大小为 $g$；最大静摩擦力近似等于滑动摩擦力，弹簧始终处于弹性限度内。求：

\begin{enumerate}
	\item
	C 在碰撞前瞬间的速度大小；
	\item
	C 与 B 碰撞过程中损失的机械能；
	\item
	从 C 与 B 相碰后到 A 开始运动的过程中，C 和 B 克服摩擦力所做的功。
\end{enumerate}

\onepicture[w=6.4cm, align=right]{figs/14.png}

%% answer:
\jdanswer{
\begin{enumerate}
	\item
	$\sqrt{v_{0}^{2}-2\mu gs_{0}}$
	\item
	$\frac{1}{4}m(v_{0}^{2}-2\mu gs_{0})$
	\item
	$\frac{2\mu^{2}m^{2}g^{2}}{k}$
\end{enumerate}
}
%% memo:
\memoanswer{
（1）小物块 C 运动至刚要与物块 B 相碰过程，根据动能定理可得

$-\mu mgs_{0}=\frac{1}{2}mv_{1}^{2}-\frac{1}{2}mv_{0}^{2}$

解得 C 在碰撞前瞬间的速度大小为

$v_{1}=\sqrt{v_{0}^{2}-2\mu gs_{0}}$

（2）物块 B、C 碰撞过程，根据动量守恒可得

$mv_{1}=2mv_{2}$

解得物块 B 与物块 C 碰后一起运动的速度大小为

$v_{2}=\frac{1}{2}\sqrt{v_{0}^{2}-2\mu gs_{0}}$

故 C 与 B 碰撞过程中损失的机械能为

$\Delta E=\frac{1}{2}mv_{1}^{2}-\frac{1}{2}\times2mv_{2}^{2}=\frac{1}{4}m(v_{0}^{2}-2\mu gs_{0})$

（3）滑板 A 刚要滑动时，对滑板 A，由受力平衡可得

$k\Delta x+2\mu mg=3\mu mg$

解得弹簧的压缩量，即滑板 A 开始运动前物块 B 和物块 C 一起运动的位移大小为

$\Delta x=\frac{\mu mg}{k}$

从 C 与 B 相碰后到 A 开始运动的过程中，C 和 B 克服摩擦力所做的功为

$W=2\mu mg\cdot\Delta x=\frac{2\mu^{2}m^{2}g^{2}}{k}$
}

\item
%% number: 15
%% paperName: 2022年福建高考第 15 题 16 分
%% typeId: 6
%% type: 计算题
%% chapter: 电磁感应
%% point: 双杆问题
%% method: 无
%% score: 16
%% degree: 350
%% duplicateId: 0
%% body:
如图~\ref{2022福建15a}，一倾角为 $\theta$ 的绝缘光滑斜面固定在水平地面上，其顶端与两根相距为 $L$ 的水平光滑平行金属导轨相连；导轨处于一竖直向下的匀强磁场中，其末端装有挡板 M、N。两根平行金属棒 G、H 垂直导轨放置，G 的中心用一不可伸长绝缘细绳通过轻质定滑轮与斜面底端的物块 A 相连；初始时刻绳子处于拉紧状态并与 G 垂直，滑轮左侧细绳与斜面平行，右侧与水平面平行。从 $t=0\Us$ 开始，H 在水平向右拉力作用下向右运动；$t=2\Us$ 时，H 与挡板 M、N 相碰后立即被锁定。G 在 $t=1\Us$ 后的速度 - 时间图线如图~\ref{2022福建15b}所示，其中 $1\sim2\Us$ 段为直线。已知：磁感应强度大小 $B=1\UT$，$L=0.2\Um$，G、H 和 A 的质量均为 $0.2\Ukg$，G、H 的电阻均为 $0.1\UO$；导轨电阻、细绳与滑轮的摩擦力均忽略不计；H 与挡板碰撞时间极短；整个运动过程 A 未与滑轮相碰，两金属棒始终与导轨垂直且接触良好：$\sin\theta=0.25$，$\cos\theta=0.97$，重力加速度大小取 $10\Umsq$，图~\ref{2022福建15b}中 $e$ 为自然常数，$\frac{4}{e}=1.47$。求：

\begin{enumerate}
	\item
	在 $1\sim2\Us$ 时间段内，棒 G 的加速度大小和细绳对 A 的拉力大小；
	\item
	$t=1.5\Us$ 时，棒 H 上拉力的瞬时功率；
	\item
	在 $2\sim3\Us$ 时间段内，棒 G 滑行的距离。
\end{enumerate}

\twopicture[labelA=2022福建15a, labelB=2022福建15b, wA=6.4cm, wB=4.4cm, align=right]{figs/15a.png}{figs/15b.png}

%% answer:
\jdanswer{
\begin{enumerate}
	\item
	$2\Umsq$，$0.9\UN$
	\item
	$16.15\UW$
	\item
	$2.53\Um$
\end{enumerate}
}
%% memo:
\memoanswer{
（1）由 $v-t$ 图像可得在 $1\sim2\Us$ 内，棒 G 做匀加速运动，其加速度为 $a=2\Umsq$ 依题意物块 A 的加速度也为 $a=2\Umsq$，由牛顿第二定律可得

$T-m_{A}g\sin\theta=m_{A}a$

解得细绳受到拉力 $T=0.9\UN$

（2）由法拉第电磁感应定律与闭合电路欧姆定律推导出“双棒”回路中的电流为

$I=\frac{BL(v_{H}-v_{G})}{R_{H}+R_{G}}$

由牛顿运动定律和安培力公式有

$BIL-T=m_{G}a$

由于在 $1\sim2\Us$ 内棒 G 做匀加速运动，回路中电流恒定为 $I=6.5\UA$，两棒速度差为 $v_{H}-v_{G}=6.5\Ums$ 保持不变，这说明两棒加速度相同且均为 $a$；

对棒 H 由牛顿第二定律可求得其受到水平向右拉力

$F=m_{H}a+BIL=1.7\UN$

由 $v-t$ 图像可知 $t=1.5\Us$ 时，棒 G 的速度为 $v_{G}=3\Ums$

此刻棒 H 的速度为 $v_{H}=9.5\Ums$

其水平向右拉力的功率 $P_{F}=Fv_{H}=16.15\UW$

（3）棒 H 停止后，回路中电流发生突变，棒 G 受到安培力大小和方向都发生变化，棒 G 是否还拉着物块 A 一起做减速运动需要通过计算判断，假设绳子立刻松弛无拉力，经过计算棒 G 加速度为

$a^{\prime}=\frac{B^{2}L^{2}v_{G}^{\prime}}{2Rm_{G}}=\frac{1^{2}\times0.2^{2}\times4}{2\times0.1\times0.2}\Umsq=4\Umsq$

物块 A 加速度为

$a^{\prime\prime}=g\sin\theta=2.5\Umsq$

说明棒 H 停止后绳子松弛，物块 A 做加速度大小为 $2.5\Umsq$ 的匀减速运动，棒 G 做加速度越来越小的减速运动；由动量定理、法拉第电磁感应定律和闭合电路欧姆定律可以求得，在 $2\sim3\Us$ 内

$\overline{B I}L\Delta t=m_{G}(v_{G2}-v_{G3})$

$\overline{I}\Delta t=\frac{BL\overline{v}}{R_{H}+R_{G}}\Delta t=\frac{BLs_{G}}{R_{H}+R_{G}}$

棒 G 滑行的距离

$s_{G}=\frac{m_{G}(v_{G2}-v_{G3})(R_{H}+R_{G})}{B^{2}L^{2}}=\left(4-\frac{4}{e}\right)\Um=2.53\Um$

这段时间内物块 A 速度始终大于棒 G 滑行速度，绳子始终松弛。
}

\end{enumerate}

\end{document}
